\documentclass[11pt, a4paper]{article}

% ===== 宏包 =====
\usepackage{ctex}
\usepackage{geometry}
\usepackage{esint}
\usepackage{enumitem}
\usepackage{amssymb}

% ===== 页面设置 =====
\geometry{left=1.8cm, right=1.8cm, top=2.5cm, bottom=2.5cm}

% ===== 文档标题 =====
\title{\Large 北京大学高等数学 A 期中试题\\[0.5em]{\normalsize 2025-2026 学年第二学期}}
\author{}
\date{}

% ===== 试题正文 =====
\setCJKmainfont{SimHei}
\begin{document}

\maketitle

 % ===== 试题内容 =====
 \begin{enumerate}
    \item (10分) 求方程 $\displaystyle \frac{\mathrm dy}{\mathrm dx} = \frac{y}{x+y^2}$ 的满足条件 $y(1)=1$ 的解.
    \item (15分) 设平面区域由曲线 $\left(2x^2+3y^2\right)^2=2x^2-3y^2$ 围成，计算积分 $\displaystyle I=\iint\limits_D \left(\sqrt{2}x+\sqrt{3}y\right)^2 \,\mathrm dx\,\mathrm dy$.
    \item (15分) 设空间区域 $\Omega$ 由曲面 $x=0$, $x=1$, $\displaystyle 2(1+x^2)=y^2+\frac{z^2}{3}$ 所围成. \\[0.6em]
        计算积分 $\displaystyle I=\iiint\limits_\Omega \left(\frac{x}{1+x^2}+\frac{y}{1+y^2}\right)\,\mathrm dx\,\mathrm dy\,\mathrm dz$.
    \item (15分) 设 $\Omega=\left\{(x,y,z) \;\Big|\; 0 \leq x-y \leq 1, 0 \leq x-z \leq 1, 0 \leq x+y+z \leq 1\right\}$. \\[0.6em]
        计算积分 $\displaystyle I=\iiint\limits_\Omega (x+y+z)\cos\left((x+y+z)^2\right) \,\mathrm dx\,\mathrm dy\,\mathrm dz$.
    \item (15分) 设 $\mathbf{F}=(P,Q,R)$，其中 $\displaystyle P=\frac{x}{(x^2+y^2+z^2)^{\frac{3}{2}}}+yz$, $\displaystyle Q=\frac{y}{(x^2+y^2+z^2)^{\frac{3}{2}}}+xz$, \\[0.6em]
        $\displaystyle R=\frac{z}{(x^2+y^2+z^2)^{\frac{3}{2}}}+xy$. 记曲面 $\Sigma^+$ 为椭球面 $\displaystyle \frac{x^2}{4}+\frac{y^2}{9}+\frac{z^2}{25}=1$ 的外侧面，\\[0.6em]
        $\mathbf{n}$ 表示 $\Sigma^+$ 的单位外法向量，计算曲面积分 $\displaystyle I=\oiint\limits_{\Sigma^+} \mathbf{F}\cdot\mathbf{n} \,\mathrm dS$.
    \item (15分) 计算 $\displaystyle I=\oint\limits_\Gamma \left(y^2-z^2\right)\,\mathrm dx+\left(z^2-x^2\right)\,\mathrm dy+\left(x^2-y^2\right)\,\mathrm dz$, \\[0.6em]
        其中 $\Gamma$ 是平面 $x+y+z=1$ 与球面 $x^2+y^2+z^2=1$ 的交线，从 $z$ 轴的正向看去逆时针方向.
    \item (15分) 
        \begin{enumerate}[label=(\arabic*)]
            \item 设 $\Omega \subset \mathbb{R}^3$ 是区域（连通开集），$f(x,y,z)$ 在 $\Omega$ 上连续. \\[0.8em]
                若在 $\Omega$ 的任意子区域 $\Omega_1$ 上均有 $\displaystyle \iiint\limits_{\Omega_1} f(x,y,z) \,\mathrm dx\,\mathrm dy\,\mathrm dz=0$，证明 $f(x,y,z)\equiv 0, (x,y,z) \in \Omega$.
            \item 设 $f(x)$ 在 $(0, +\infty)$ 具有连续的一阶导数，且 $\lim\limits_{x \to 0^+} f(x)$ 存在. 又对 $\mathbb{R}^3_{x>0}=\left\{(x,y,z) \;\Big|\; x>0\right\}$ \\[0.6em]
                中的任意光滑简单闭曲线 $S$（外侧）都有 $\displaystyle \oiint\limits_S xf(x)\,\mathrm dy\,\mathrm dz-xyf(x)\,\mathrm dz\,\mathrm dx+z\mathrm e^{2x}\,\mathrm dx\,\mathrm dy=0$. \\[0.6em]
                求 $f(x)$ 的表达式.
        \end{enumerate}
 \end{enumerate}

\end{document}